\section{数据中心经济学：从 grid MW 到 useful tokens}

有了芯片并不等于拥有可用算力。电网输入要经过供配电、UPS、cooling、network、storage、CPU 和 accelerator，最终只有一部分变成满足质量与延迟目标的 useful work。本节把讲者的 gigawatt 和 facility-cost 讨论转成资源会计框架，不复述未经审计的绝对数字。

\subsection{PUE、利用率与 workload efficiency}

\term{PUE}（Power Usage Effectiveness）定义为
\[
\mathrm{PUE}=\frac{P_{facility}}{P_{IT}},
\]
其中 $P_{facility}$ 是数据中心总功率，$P_{IT}$ 是服务器、network 和 storage 等 IT load。PUE 越接近 1，facility overhead 越低；但它不说明 accelerator 是否被充分利用，也不说明模型是否以最低成本满足 SLO。

\lecturefigure{05-power-budget.png}{Facility power 经过 overhead 与 IT load 后，才产生 useful accelerator work。}{本地字幕 10:30--14:50；概念重绘。}

\begin{knowledgebox}{读图：三个效率指标缺一不可}
\textbf{Facility efficiency} 看 PUE；\textbf{hardware utilization} 看 accelerator busy time、memory 和 network stall；\textbf{workload efficiency} 看 tokens、samples 或 training progress 相对于资源。低 PUE 的空闲集群仍然浪费，高利用率但生成失败结果的系统也没有价值。
\end{knowledgebox}

\begin{importantbox}{Useful compute per grid watt}
可用一个诊断式表达
\[
\eta_{grid}=\frac{U\cdot \eta_{workload}}{\mathrm{PUE}},
\]
其中 $U$ 是 IT 资源有效利用率，$\eta_{workload}$ 是单位 IT power 产生的合格工作量。提高 $U$ 可能依靠 batching、scheduling 与 demand aggregation；提高 $\eta_{workload}$ 可能依靠更好的 kernel、model routing、quantization 和 cache。
\end{importantbox}

\subsection{规划时钟：电力、建筑、芯片、模型并不同步}

国家级项目最危险的并非“预测不准”，而是把不同寿命的决策锁在一起。grid 与 site 可能按十年规划；facility shell 需要多年建设；accelerator 与 rack density 快速变化；model demand 又受算法和产品影响。读者需要把不可逆的土建与可替换的 compute module 分层。

\lecturefigure{06-planning-clocks.png}{国家 AI 基础设施同时运行在四个互相冲突的时间尺度上。}{本地字幕 14:45--16:10；概念重绘。}

\begin{knowledgebox}{读图：用 modularity 管理 regret}
电力与土地层保留 headroom；facility 支持更高密度 cooling 和 network retrofit；rack 与 accelerator 尽量模块化；software 保持模型和硬件可移植。Modularity 会增加初始成本，却能降低“数据中心上线时已经不适合新一代 rack”的路径依赖。
\end{knowledgebox}

\begin{warningbox}{计划容量不等于在线容量}
现场提到的 1.5--2 GW、建设周期和单位成本属于课堂口径。正式评估必须区分 announced、contracted、under construction、energized、installed、available 与 utilized capacity；否则一个数字会混合土地、电力、building shell 和可服务 token 的不同阶段。
\end{warningbox}

\teachervoice{讲者的工程直觉是“尽量把电送到 accelerator”。更完整的讲义版本应再加一句：accelerator 只有被可靠 workload 使用时才是生产资产；capacity planning 必须同时建设 scheduler、software、人才、模型 onboarding 与 demand pipeline。}

\subsection{本章小结}

数据中心经济学必须从 grid watt 一直核算到 useful outcome。PUE、utilization、workload efficiency 和 planning clocks 共同决定总成本，单独展示任何一个指标都可能误导。

